\subsection{Platform and Robot}
Experiments use the pinned Arena ROS2 Humble runtime with the Gazebo fallback,
Nav2 NavFn/DWB, and a Jackal model.  The map spans
$31.28\times24.03$ m; evaluation uses a 0.36-m robot radius and 0.35-m
pedestrian radius.  A planar LiDAR is resampled to the 180 bins consumed by the
recovery policy.  Localized and physical goal tolerances are 0.25 and 0.30 m,
respectively.  The deployment interface contains only LiDAR, path, velocity,
command, progress, planner-status, and rule-detector histories; simulator poses
are isolated for actor execution, expert labeling, and evaluation.

Pedestrians are LiDAR-visible cylindrical actors following deterministic,
one-shot routes at 5 Hz, with family-dependent speeds spanning 0.06--0.55 m/s.
Their swept-step guard holds a candidate motion that would reduce robot
clearance inside the configured 0.90 m soft-yield region; the 0.73 m hard guard
prevents a step from crossing the robot footprint.  This protocol makes
pedestrian motion replayable without treating scripted actor state as a
deployment observation.  It is an interaction stress model rather than a
model of human intent.

\IfFileExists{figures/runtime_gazebo_doorway_bottleneck_medium.png}{%
\begin{figure}[t]
  \centering
  \includegraphics[width=\columnwidth]{runtime_gazebo_doorway_bottleneck_medium.png}
  \caption{Actual Arena/Gazebo execution of the doorway interaction with
  Jackal and LiDAR-visible pedestrians.  The image is qualitative runtime
  evidence; all numerical results are computed from paired episode logs.}
  \label{fig:runtime-gazebo}
\end{figure}
}{}

\subsection{Moderate-v5 Scenarios and Splits}
Figure~\ref{fig:scenarios} summarizes the head-on corridor, doorway bottleneck, crossing flow,
blind corner, group blocking, overtaking, opposite streams, and temporary
blockage.  Low, medium, and high densities contain one, two, and four
pedestrians, respectively.  All splits intentionally share the
\texttt{map\_empty} map and the eight family templates; train, validation, and
test instead use disjoint seed blocks, scenario IDs, and compiled physical
realizations.  Thus the protocol tests held-out interactions, not unseen-map or
unseen-template generalization.

The validation manifest contains three repetitions for each of the 24
family--density cells: 72 conditions per method and 360 logical
method--episodes across the five methods.  These runs are reserved for code,
checkpoint, and configuration selection.  The held-out test plan contains five
repetitions per cell: 120 conditions per method and 600 logical
method--episodes total.  Every method receives the same condition keys within
each split.

Moderate-v5 was constructed from moderate-v4 validation evidence, which exposed
static-geometry overlaps and interaction lanes narrower than the combined
robot--pedestrian collision radii.  V5 repairs those construction defects while
retaining the eight families, density levels, footprint, planner, and actor
protocol.  The full v5 validation manifest is the only split authorized for
checking reference-planner difficulty and interaction coverage.  No v5 test
episode is executed or inspected before the code, configuration, scenario
hashes, and checkpoints are frozen.

All five methods run the same condition keys, robot starts, goals, pedestrian
routes, seeds, 240-s timeout, and planner profile.  Final generation fails if a
method is missing a condition, if scenario metadata differs across a pair, or
if the observed condition count is not the frozen value.

\subsection{Baselines and Learned Method}
The closed-loop comparison includes: (i) Base DWB without a recovery subtree;
(ii) the Standard Nav2 recovery behavior; (iii) a rule-triggered Heuristic
policy with the same state machine, planning mask, and safety supervisor as the
learned hierarchy; (iv) Uniform BC using the planning-expert demonstrations;
and (v) PGRR using the validation-selected DAgger checkpoint.  The first four
are complete paired baselines rather than a high-density subset.  The
privileged short-horizon expert is used as a planning reference for labels and
diagnostics only; it is not a deployment baseline, and no optimality claim is
made.

Training and model selection use only train and validation scenarios.  The
two-round DAgger workflow is preserved, but the second-round candidate did not
pass validation selection.  The deployed checkpoint instead combines the
accepted DAgger aggregate with a train-only head-on coverage shard and is chosen
by validation loss and closed-loop validation behavior.  Margin weighting
remains a reported negative offline ablation.  PPO and a learned failure
detector are outside the evaluated method and are not claimed as completed
contributions.

\subsection{Metrics and Statistical Protocol}
Every stream records physical robot and pedestrian geometry, commands, planner
and recovery state, failure scores, LiDAR, and paths.  We report goal reaching,
collision, timeout, planner failure, SPL, path length, successful-run navigation
time, minimum human-centre distance, personal-space violation, discomfort time,
emergency stops, triggers, recovery success, recovery duration, and
intervention ratio.  \texttt{SIMULATOR\_FAILURE} and
\texttt{INVALID\_RESET} attempts remain in the artifacts and may be excluded
from algorithm rates according to the declared protocol; collision, timeout,
and planner failure are never merged or removed.

A recovery success is a logged REJOIN-to-NORMAL transition.  The state-machine
guard requires the original task goal to be active again, failure score below
$\tau_{off}$, and valid goal progress.  Reaching the final goal is a separate
navigation outcome, and the recovery metric does not add an unimplemented
post-transition protection window.

PGRR is paired separately with each baseline on identical condition keys.
Binary outcomes use exact McNemar tests and episode-pair bootstrap confidence
intervals.  Continuous paired differences use two-sided Wilcoxon signed-rank
tests, paired bootstrap intervals, and rank-biserial effect sizes.  One global
Holm correction covers all predeclared comparator--endpoint tests.  Confidence
intervals use 10,000 paired bootstrap resamples with recorded seed 20260804.
The tables report effect sizes and intervals regardless of statistical
significance.

\subsection{Execution and Artifact Integrity}
Six concurrent workers use isolated ROS domain IDs and simulator partitions.  A
startup handshake holds actors until navigation, authoritative pose feedback,
and logging are active.  The runner refuses overwrite, records checkpoint and
scenario hashes, and retains every physical attempt.  A logical task permits at
most three physical attempts: the initial attempt and up to two retries, each
only after an explicitly classified \texttt{SIMULATOR\_FAILURE} or
\texttt{INVALID\_RESET}.  Collision, timeout, and planner failure are accepted
algorithm outcomes and are never rerun because they are unfavorable.  Exactly
one accepted terminal result is required per logical task.  The paper pipeline
first collects the locked run, then validates the complete five-method
condition set, computes statistics, and generates all result tables and figures
from Parquet, CSV, and JSON artifacts.
